% Tagged PDF/UA-2 fixture for test/pdf-a11y-test.el.
%
% Rebuild a11y.pdf from this directory with (run twice for the
% outline and structure tree to settle):
%
%   lualatex a11y.tex && lualatex a11y.tex
%
% then delete the .aux and .log files.  The images come from the mwe
% package in TeX Live.  The figures cover good alt text, alt text
% that is only a file name, and no alt text in the source.  For the
% last one the LaTeX tagging code itself writes the file name
% (example-image-b.pdf) as /Alt, so neither figure is flagged by
% veraPDF, which is what the report is for.
%
% LaTeX never writes a figure without /Alt, so a11y-noalt.pdf is
% derived from a11y.pdf with qpdf: the file-name alt becomes empty
% and the second one is removed.
%
%   qpdf --qdf --object-streams=disable a11y.pdf tmp.pdf
%   LC_ALL=C sed \
%     -e '/\/Alt <feff006500780061006d0070006c0065002d0069006d006100670065005f/s/<.*>/()/' \
%     -e '/\/Alt <feff006500780061006d0070006c0065002d0069006d006100670065002d0062/d' \
%     tmp.pdf | fix-qdf > tmp2.pdf
%   qpdf tmp2.pdf a11y-noalt.pdf && rm tmp.pdf tmp2.pdf
\DocumentMetadata{lang=en, pdfstandard=ua-2, testphase=latest}
\documentclass{article}
\usepackage{graphicx}
\title{Accessibility Fixture}
\author{pdf-tools}
\date{}
\begin{document}
\maketitle
\section{Introduction}
This document exercises the accessibility report.

\begin{figure}[h]
  \centering
  \includegraphics[width=3cm,alt={A grey placeholder rectangle with the word Image}]{example-image}
  \caption{A figure with good alternative text.}
\end{figure}

\subsection{Bad alternative text}
The next figure is described only by its file name.

\includegraphics[width=3cm,alt={example-image_a.png}]{example-image-a}

The next figure has no alternative text at all.

\includegraphics[width=3cm]{example-image-b}

\newpage
\section{Second page}
Text on the second page.
\end{document}
